\subsection*{MT-M31. Local Euler differences: actual range, Hilbert closure,
and the Schanuel dependency}

The distinction in this result is between the range of a bounded
operator and the closure of that range.  It occurs before any
GRH or positivity question.

\paragraph{The precise source assertion.}
Masullo, Proposition 11.2.0.1(2), original PDF pp.~484--485,
asserts that the extension of $\mathrm{PV}_{S_0}$ to the
zero-line Hilbert completion has closed range when
$s\geq 2|S_0|$.  Equation (11.2.0.3), p.~484, identifies
the closure of the smooth image with that actual operator range.
The proof on p.~484 gives, for $S_0=\{v\}$, the exact formula
\[
  P_v\Phi_F(f)=\Phi_F(P_vf),\qquad
  P_vf(x)=f(x)-f(\pi_v^{-1}x).
  \tag{M31.1}
\]
There is no additional norm-character multiplier in this formula.
On p.~485 the proof identifies the range with an intersection
of kernels of continuous coefficient functionals.  The argument
below proves that this identifies the \emph{closure}, and that
the actual range is not closed even for one prime over $\mathbf Q$.
The original images of both pages have been read, including
the multiplier and the range/closure equation.

\paragraph{General analytic formulation.}
For $s\geq0$ put
\[
 H_s=L^2(\mathbf R,w_s(x)\,dx),\qquad
 w_s(x)=(1+x^2)^{s/2},\qquad
 (U_hf)(x)=f(x-h).
 \tag{M31.2}
\]
This is precisely the radial form of the source norm
(9.1.2.3), original p.~399, with the $a=0$ completion
of (10.1.1.3), p.~434; see MT-M24--M25 for the actual
theta-image identification.  Fix a nonempty finite list
of $h_v>0$ and $a_v\in\mathbf C$ with $|a_v|=1$, and set
\[
 P=\prod_v(1-a_vU_{h_v}),\qquad
 D(z)=\prod_v(1-a_ve^{h_vz}).
 \tag{M31.3}
\]
Repeated factors are allowed.  Let $\mathcal T\subset\mathbf R$
be the locally finite set of $t$ with $D(it)=0$, and let
$m_t$ be its multiplicity.  Define
\[
 q_s=\max\{0,\lceil(s-1)/2\rceil\},\qquad
 J_{t,j}(f)=\int_{\mathbf R}f(x)x^je^{itx}\,dx.
 \tag{M31.4}
\]
The following assertions hold without any conjectural input:
\begin{enumerate}
\item $P$ is a bounded injective operator on $H_s$, and its
actual range is not closed, for every $s\geq0$.
\item Its exact closed range replacement is
\[
 \overline{P H_s}^{\,H_s}
 =\bigcap_{t\in\mathcal T}\ \bigcap_{0\leq j<\min(m_t,q_s)}
            \ker J_{t,j}.
 \tag{M31.5}
\]
Every displayed functional is continuous.  Thus for $0\leq s\leq1$
the actual range is dense and proper; for $s>1$ its closure
is proper and the range is a proper dense subspace of that closure.
\item If $\mathcal A$ is any dense linear subspace of $H_s$, then
\[
 \overline{P\mathcal A}^{\,H_s}
       =\overline{PH_s}^{\,H_s}.
 \tag{M31.6}
\]
In particular the corrected quotient is the Hilbert quotient by
(M31.5).  It is not the non-Hausdorff topological quotient
$H_s/PH_s$.
\end{enumerate}

\paragraph{Boundedness and injectivity.}
The inequality
\[
 1+(x+y)^2\leq2(1+x^2)(1+y^2)
 \tag{M31.7}
\]
shows that $\|U_h\|_{H_s\to H_s}\leq
2^{s/4}(1+h^2)^{s/4}$.  It gives the analogous estimate
on $H_{-s}=L^2(w_s^{-1}dx)$ as well.  Each factor and
their product are bounded.  Since $w_s\geq1$, $H_s$
embeds in ordinary $L^2(\mathbf R)$.  The ordinary Fourier
transform of $Pf$ is the transform of $f$ multiplied by
\[
            \prod_v(1-a_ve^{-ih_v\xi}).
\]
Its real zeros form a locally finite set, of Lebesgue measure
zero.  Consequently $Pf=0$ implies $f=0$ in $L^2$,
hence in $H_s$.

\paragraph{A quantitative failure of closed range.}
Choose one root $t_0$ of one of the factors of $D(it)$,
and choose nonzero $\psi\in C_c^\infty((1,2))$.
For $R\geq1$ put
\[
             f_R(x)=e^{-it_0x}\psi(x/R).
 \tag{M31.8}
\]
Changing variables gives constants $c,C>0$, independent of $R$,
such that
\[
      cR^{s+1}\leq\|f_R\|_{H_s}^2\leq CR^{s+1}.
 \tag{M31.9}
\]
For a factor with $a_ve^{it_0h_v}=1$,
\[
 (1-a_vU_{h_v})f_R(x)
   =e^{-it_0x}\{\psi(x/R)-\psi((x-h_v)/R)\}.
 \tag{M31.10}
\]
The mean-value bound, with support contained in an interval
of length $O(R)$ and $|x|=O(R)$, gives
\[
 \|(1-a_vU_{h_v})f_R\|_{H_s}
      \leq C_vR^{(s-1)/2}.
\]
The remaining finitely many factors are bounded independently
of $R$.  Since they commute,
\[
       \frac{\|Pf_R\|_{H_s}}{\|f_R\|_{H_s}}=O(R^{-1}).
 \tag{M31.11}
\]
An injective bounded operator between Hilbert spaces with
closed range has bounded inverse onto that range, by the
bounded-inverse theorem, and hence is bounded below.  Equation
(M31.11) contradicts this.  It proves the first assertion.

\paragraph{The exact closure.}
Use the continuous bilinear duality
\[
       H_s\times H_{-s}\longrightarrow\mathbf C,
       \qquad(f,g)\longmapsto\int f(x)g(x)\,dx.
 \tag{M31.12}
\]
Complex conjugation converts it to the usual Hilbert duality;
the bilinear version keeps the Mellin convention explicit.
The annihilator $\mathcal N$ of $PH_s$ consists precisely of
the $g\in H_{-s}$ satisfying
\[
          P^{\mathrm t}g=
             \prod_v(1-a_vU_{-h_v})g=0.
 \tag{M31.13}
\]
Each $g\in H_{-s}$ is a tempered distribution: Cauchy--Schwarz
bounds its pairing with a Schwartz test by
$\|g\|_{H_{-s}}\|\phi\|_{H_s}$.  With Fourier convention
$\widehat g(\xi)=\int g(x)e^{-i\xi x}dx$, (M31.13) is
\[
                         D(i\xi)\widehat g(\xi)=0.
 \tag{M31.14}
\]
It follows that $\widehat g$ is supported on $\mathcal T$.
Near an isolated $t\in\mathcal T$, a distribution killed
by the smooth function $D(i\xi)$ is a linear combination
of $\delta_t,\ldots,\delta_t^{(m_t-1)}$.  To verify the
order bound, divide by the nonvanishing smooth factor in
$D(i\xi)=(\xi-t)^{m_t}b(\xi)$ and use
$(\xi-t)\delta_t^{(j)}=-j\delta_t^{(j-1)}$.
The classical point-support distribution theorem,
\emph{Expository Notes on Distribution Theory}, AMSC/MATH~673,
Fall~2018, Theorem~2.1, p.~4 (University of Maryland),
then gives exactly this finite list.

Here localization is bounded in the required Hilbert norm.
For any Schwartz kernel $k$, (M31.7) and Minkowski give
\[
 \|k*g\|_{H_{-s}}
 \leq C_s\left(\int |k(y)|(1+|y|)^{s/2}dy\right)
                 \|g\|_{H_{-s}}.
 \tag{M31.15}
\]
Multiplying $\widehat g$ by a compactly supported smooth
function supported near one $t$ is therefore bounded on
$H_{-s}$.  The resulting inverse transform is a polynomial
times $e^{itx}$.  A nonzero polynomial of degree $j$
times this exponential belongs to $H_{-s}$ exactly when
\[
                              2j<s-1.
 \tag{M31.16}
\]
Indeed its squared absolute value is comparable to $|x|^{2j}$
for all sufficiently large $|x|$; the weighted integral
then has tail comparable to $\int_1^\infty x^{2j-s}dx$.
In particular equality in (M31.16) is not allowed.
This argument isolates each frequency first, so cancellation
among distinct frequencies cannot evade the cutoff.

For completeness the finite exponential-polynomial sums are
dense in $\mathcal N$ in its actual $H_{-s}$ norm.
Choose $\eta\in C_c^\infty(\mathbf R)$ equal to one near zero,
and put $\eta_R(\xi)=\eta(\xi/R)$, $R\geq1$.
Its inverse Fourier kernel has the form $k_R(x)=Rk(Rx)$,
with integral one.  Equation (M31.15) gives uniformly bounded
operators because
\[
 \int |k_R(y)|(1+|y|)^{s/2}dy
   =\int |k(u)|(1+|u|/R)^{s/2}du
   \leq\int |k(u)|(1+|u|)^{s/2}du.
\]
Translations are strongly continuous on $H_{-s}$, first on
$C_c^\infty$ and then by density and (M31.7).
Writing $k_R*g-g$ as the integral of
$k(u)(U_{u/R}g-g)$ proves convergence to zero in $H_{-s}$
by dominated convergence, using the same polynomial bound.
Each $\eta_R\widehat g$ has finite support in $\mathcal T$,
so its inverse transform is a finite sum of the terms
$x^je^{itx}$ with $j<m_t$ and (M31.16).
Conversely all these terms satisfy (M31.13), and their
closed span therefore equals $\mathcal N$.

Taking the annihilator once more proves (M31.5).
Cauchy--Schwarz proves continuity of every $J_{t,j}$ in that
formula using (M31.16).  For $s>1$, at least one root exists
and $J_{t,0}$ is a nonzero functional, so the closure is proper.
For $s\leq1$ no term survives and the closure is all of $H_s$.
Nonclosedness gives all remaining strictness assertions.
Finally boundedness of $P$ and density of $\mathcal A$ prove
(M31.6) directly by approximating each vector before applying $P$.

\paragraph{Application to the literal adelic assertion.}
It is enough to take $F=\mathbf Q$, $S_0=\{p\}$, and
the trivial character of the compact norm-one group.
Choose a norm splitting $C_{\mathbf Q}\simeq
C_{\mathbf Q}^1\times\mathbf R$ with radial coordinate
$x=\log|g|$, and normalize the compact Haar measure to one.
On this reducing sector (M31.1) is
$f(x)\mapsto f(x)-f(x+\log p)$.
Replacing $x$ by $-x$ gives exactly $1-U_{\log p}$ on
(M31.2).  Reflection is isometric because $w_s$ is even.

The ambient operator on the full adelic Hilbert space is
injective as well.  If $\Phi(g)=\Phi(\pi_p^{-1}g)$
almost everywhere, decompose $C_{\mathbf Q}$ into the
disjoint norm shells
$\{k\log p\leq\log|g|<(k+1)\log p\}$.
Translation by $\pi_p^{-1}$ maps successive shells onto
one another and preserves Haar measure.  The unweighted
$L^2$ mass is consequently equal on all shells.  Since
$H_s\subset L^2$ for $s\geq0$, it must be zero.
Thus the restriction to any closed invariant Hilbert subspace
is also injective.  If its range were closed, that restriction
would be bounded below on every sector it contains.

The actual zero-line completion of the trivial-character
theta image contains the whole radial $H_s$ sector.
This follows from MT-M24 and the exact MT-M25 kernel formula
at $a=0$, since the pole-removed completed zeta function has
no nontrivial zero on $\Re z=0$.  The latter is the classical
zero-free boundary theorem and functional equation, with no
RH premise.  Thus the normalized vectors from (M31.8)
belong to the actual completion, and (M31.11) contradicts
Proposition 11.2.0.1(2), including its stated range $s\geq2$.
This proves a counterexample in the specified space, not only
in an unrelated radial model.

By (M31.6), the left side of source (11.2.0.3) is precisely
the closure in (M31.5); its right side, read as the actual
extended-operator range, is a proper dense subspace of it.
Replacing the right side by its closure is the necessary
correction.  No conclusion about the truth of $M_H$ follows
from this counterexample.

\paragraph{What the local-pole input actually assumes.}
The following is a targeted dependency trace, not an audit
of the entire draft.
\begin{enumerate}
\item Theorem 7.1.3.2(1), original pp.~324--325, says that
two unramified finite places of distinct residue
characteristics have at most one common Euler pole.
If $s,s'$ were distinct common poles, then
$(s-s')\log q_v$ and $(s-s')\log q_w$ would both be
nonzero integral multiples of $2\pi i$, forcing
$\log q_v/\log q_w\in\mathbf Q$, impossible for powers
of distinct rational primes.  This part is unconditional.
\item Part (3b) of that theorem explicitly assumes Schanuel
for its sharper coincidence exclusion for a character
which is an algebraic character times a norm power.
The proof on p.~325 uses the four-exponentials consequence
of Schanuel.  Its weaker three-residue-characteristic
exclusion uses the unconditional six-exponentials theorem.
The source does not present the four-exponentials assertion
as an unconditional input in that argument.
\item For one rational prime $p$, the pole multiplicity of
$\prod_{v\mid p}L_v(s,\chi_v)$ is at most the number
of factors, hence at most $[F:\mathbf Q]$.
Each unramified Euler factor has simple poles, and ramified
factors contribute none.  This elementary bound does not
use Schanuel or coincidence control between different primes.
Definition 7.3.1.4, p.~343, and Theorem 7.3.1.6,
pp.~344--350, organize the construction by precisely these
$p$-primary multiplicities.  The supplied proof of the latter
was read as a dependency map; its invoked structure and
interpolation theorems are not certified here.
\item The source's zero-line coefficient norm is
$\int_{\mathbf R}|x|^{2j}(1+x^2)^{-s/2}dx$,
displayed in Proposition 11.2.0.4, p.~486.
The strict condition is (M31.16), not a conjectural
arithmetic statement.  If $j\leq [F:\mathbf Q]-1$
and $s\geq2[F:\mathbf Q]$, it holds strictly.
Thus the stated working degree $2[F:\mathbf Q]$
retains every $p$-primary jet without Schanuel.
The source's general boundary cutoff (11.2.0.6), p.~485,
has the same equality problem already corrected by MT-M25;
it does not affect this sufficient even-integer degree.
\item Theorem 11.2.2.1 and Remark 11.2.2.2,
pp.~490--491, use $s\geq2[F:\mathbf Q]$ and
$N(\varepsilon_p)\leq[F:\mathbf Q]$ to assert a
trace with full local multiplicity.  The statement and its
complete supplied proof contain no Schanuel premise.
Their Hilbert quotient and coefficient construction instead
depend on the actual/closed-range identification above.
The correct local quotient can be constructed using
(M31.5); the further circle model and local trace are
treated separately in MT-M32.
\end{enumerate}

\paragraph{The precise downstream gate.}
The dependency is explicit in the text, rather than inferred
from proximity of the statements.  Theorem 11.2.0.2,
pp.~485--486, uses the space asserted closed in
(11.2.0.3); the definition of the orbital Hilbert quotient
at pp.~489--490 and Theorem 11.2.2.1 then use this space.
Theorem 11.3.2.1, pp.~495--496, identifies a localized
geometric kernel with the local orbital action and invokes
the finite-rank trace approximation (11.2.2.12) from
Theorem 11.2.2.1.  Finally the proof of Theorem 12.5.1.1,
p.~520, explicitly invokes Theorem 11.3.2.1 for the
contributions away from the chosen finite set of places,
and Theorem 12.4.1.1 for its remaining local contribution.
The supplied statements and complete displayed proof of
11.3.2.1 were read for this trace.  The full proof of
12.4.1.1 is not claimed audited here.

Consequently the local closure correction preserves a
meaningful local orbital quotient, but does not by itself
prove the geometric-kernel identification in 11.3.2.1
or the narrow-strip trace identity of 12.5.1.1.
The Schanuel refinement belongs to a different branch:
the introductory Expectation 1.7.1, original pp.~37--38,
labels the sharper zero-width multiplicity consequences
as expectations not fully worked out.  No hidden Schanuel
assumption has been established in the positive-width
geometric trace route by this targeted dependency reading.
Its remaining issue is the stated analytic comparison,
not the elementary $p$-primary pole bound.
